% --- Archon physics formalization source begin ---
% archon:physics
% archon:covers PhyXMiniProblems/problem_phyx_mini_0977.lean
% archon:source-report reports/phyx_mini/problem_phyx_mini_0977.source.json

\chapter{Physics problem phyx\_mini\_0977}
\label{ch:PhyXMiniProblems_problem_phyx_mini_0977}

\paragraph{Problem source.}
\#\# Physical scenario

A rectangular loop with width \textbackslash{}( L \textbackslash{}) and a slidewire with mass \textbackslash{}( m \textbackslash{}) are as shown. A uniform magnetic field \textbackslash{}( \textbackslash{}vec\{B\} \textbackslash{}) is directed perpendicular to the plane of the loop into the plane of the figure. The slidewire is given an initial speed of \textbackslash{}( v\_0 \textbackslash{}) and then released. There is no friction between the slidewire and the loop, and the resistance of the loop is negligible in comparison to the resistance \textbackslash{}( R \textbackslash{}) of the slidewire.

\#\# Question

Show that the distance \textbackslash{}( x \textbackslash{}) that the wire moves before coming to rest is \textbackslash{}( x = \textbackslash{}frac\{mv\_0 R\}\{L\textasciicircum{}2 B\textasciicircum{}2\} \textbackslash{}).

\#\# Answer choices

- A: A: \textbackslash{}frac\{mv\_0\}\{RB\textasciicircum{}2 L\textasciicircum{}2\}

- B: B: \textbackslash{}frac\{mv\_0 R\}\{B\textasciicircum{}2 L\textasciicircum{}2\}

- C: C: \textbackslash{}frac\{mv\_0 R\}\{B\textasciicircum{}2 L\}

- D: D: \textbackslash{}frac\{mv\_0 R\}\{B L\textasciicircum{}2\}

\#\# Figure caption (auxiliary; use the image as primary evidence)

The image depicts a U-shaped conductor with a sliding rod across its arms, forming a closed loop. The region inside the loop is filled with a uniform magnetic field represented by blue "X"s, indicating that the magnetic field is directed into the page. 

- There is a green arrow labeled \textbackslash{}( v \textbackslash{}) along the sliding rod, pointing right, indicating the velocity of the rod.
- The magnetic field is represented by the symbol \textbackslash{}( \textbackslash{}vec\{B\} \textbackslash{}).
- The length of the rod within the field is labeled as \textbackslash{}( L \textbackslash{}).
- The background is uniformly filled with the "X" pattern, emphasizing the consistent direction of the magnetic field across the entire area inside the loop.

This setup is commonly used to illustrate electromagnetic induction principles.

\#\# Dataset metadata

Electromagnetic Induction; Physical Model Grounding Reasoning, Multi-Formula Reasoning

\paragraph{Recorded answer/context.}
B: B: \textbackslash{}frac\{mv\_0 R\}\{B\textasciicircum{}2 L\textasciicircum{}2\}

\paragraph{Figure/image path.}
/root/proposal\_for\_physic/science-mango/phyx\_mini\_run/phyx\_data/test\_image/977.png

\paragraph{Formalization target.}
create a compiling Lean file with sorry bodies at `PhyXMiniProblems/problem\_phyx\_mini\_0977.lean`. The Lean declarations must preserve the physical quantities, dimensions or dimensional roles, figure labels, governing-law hypotheses, and final relation expressed by this problem.
Use Mathlib/Physlib names found through LeanExplore where available. If a physics API is missing, introduce faithful local abstractions rather than scalar placeholder aliases.

\begin{theorem}[Physics formalization target]
\label{thm:physics:phyx_mini_0977:target}
\lean{PhyXMiniProblems.ProblemPhyXMini0977.problem_phyx_mini_0977}
\uses{def:physics:phyx-mini-0977:phyxminiproblems-problemphyxmini0977-lengthinmeters, def:physics:phyx-mini-0977:phyxminiproblems-problemphyxmini0977-massinkilograms, def:physics:phyx-mini-0977:phyxminiproblems-problemphyxmini0977-speedmagnitudeinmeterspersecond, def:physics:phyx-mini-0977:phyxminiproblems-problemphyxmini0977-magneticfluxdensityinteslas, def:physics:phyx-mini-0977:phyxminiproblems-problemphyxmini0977-resistanceinohms, def:physics:phyx-mini-0977:phyxminiproblems-problemphyxmini0977-slidewirebrakingsetup, def:physics:phyx-mini-0977:phyxminiproblems-problemphyxmini0977-matchesidealslidewirescenario, def:physics:phyx-mini-0977:phyxminiproblems-problemphyxmini0977-matchesprimaryslidewirefigure, def:physics:phyx-mini-0977:phyxminiproblems-problemphyxmini0977-hasphysicalslidewireparameters, def:physics:phyx-mini-0977:phyxminiproblems-problemphyxmini0977-modelsuniformintopagemagneticfield, def:physics:phyx-mini-0977:phyxminiproblems-problemphyxmini0977-satisfiesslidewirekinematics, def:physics:phyx-mini-0977:phyxminiproblems-problemphyxmini0977-satisfiesmotionalemflaw, def:physics:phyx-mini-0977:phyxminiproblems-problemphyxmini0977-satisfiesclosedloopohmslaw, def:physics:phyx-mini-0977:phyxminiproblems-problemphyxmini0977-satisfiesmagneticbrakingforcelaw, def:physics:phyx-mini-0977:phyxminiproblems-problemphyxmini0977-satisfiesnewtonsecondlaw, def:physics:phyx-mini-0977:phyxminiproblems-problemphyxmini0977-hasasymptoticstoppingdistance}
The assigned autoformalize agent should translate this physics problem into Lean declarations in the covered file, with theorem and lemma proof bodies written as `by sorry`.
\end{theorem}
\begin{proof}
This is an autoformalization task, not a proof task. Produce statements that can later be proved by the physics prover without weakening the physical model.
\end{proof}

% --- Archon declaration topology begin ---
\section{Declaration topology}
\begin{definition}[speed Dimension]
  \label{def:physics:phyx-mini-0977:phyxminiproblems-problemphyxmini0977-speeddimension}
  \lean{PhyXMiniProblems.ProblemPhyXMini0977.speedDimension}
  Speed has physical dimension L T⁻¹
\end{definition}
\begin{proof}
  This is the declared model definition of speed Dimension.
\end{proof}
\begin{definition}[acceleration Dimension]
  \label{def:physics:phyx-mini-0977:phyxminiproblems-problemphyxmini0977-accelerationdimension}
  \lean{PhyXMiniProblems.ProblemPhyXMini0977.accelerationDimension}
  Acceleration has physical dimension L T⁻²
\end{definition}
\begin{proof}
  This is the declared model definition of acceleration Dimension.
\end{proof}
\begin{definition}[magnetic Flux Density Dimension]
  \label{def:physics:phyx-mini-0977:phyxminiproblems-problemphyxmini0977-magneticfluxdensitydimension}
  \lean{PhyXMiniProblems.ProblemPhyXMini0977.magneticFluxDensityDimension}
  Magnetic flux density has physical dimension M T⁻¹ C⁻¹ (tesla)
\end{definition}
\begin{proof}
  This is the declared model definition of magnetic Flux Density Dimension.
\end{proof}
\begin{definition}[electromotive Force Dimension]
  \label{def:physics:phyx-mini-0977:phyxminiproblems-problemphyxmini0977-electromotiveforcedimension}
  \lean{PhyXMiniProblems.ProblemPhyXMini0977.electromotiveForceDimension}
  Electromotive force has dimension M L² T⁻² C⁻¹ (volt)
\end{definition}
\begin{proof}
  This is the declared model definition of electromotive Force Dimension.
\end{proof}
\begin{definition}[electric Current Dimension]
  \label{def:physics:phyx-mini-0977:phyxminiproblems-problemphyxmini0977-electriccurrentdimension}
  \lean{PhyXMiniProblems.ProblemPhyXMini0977.electricCurrentDimension}
  Electric current has dimension C T⁻¹ (ampere)
\end{definition}
\begin{proof}
  This is the declared model definition of electric Current Dimension.
\end{proof}
\begin{definition}[electrical Resistance Dimension]
  \label{def:physics:phyx-mini-0977:phyxminiproblems-problemphyxmini0977-electricalresistancedimension}
  \lean{PhyXMiniProblems.ProblemPhyXMini0977.electricalResistanceDimension}
  Electrical resistance has dimension M L² T⁻¹ C⁻² (ohm)
\end{definition}
\begin{proof}
  This is the declared model definition of electrical Resistance Dimension.
\end{proof}
\begin{definition}[force Dimension]
  \label{def:physics:phyx-mini-0977:phyxminiproblems-problemphyxmini0977-forcedimension}
  \lean{PhyXMiniProblems.ProblemPhyXMini0977.forceDimension}
  Force has physical dimension M L T⁻² (newton)
\end{definition}
\begin{proof}
  This is the declared model definition of force Dimension.
\end{proof}
\begin{definition}[Length Quantity]
  \label{def:physics:phyx-mini-0977:phyxminiproblems-problemphyxmini0977-lengthquantity}
  \lean{PhyXMiniProblems.ProblemPhyXMini0977.LengthQuantity}
  A nonnegative, unit-independent physical length
\end{definition}
\begin{proof}
  This is the declared model definition of Length Quantity.
\end{proof}
\begin{definition}[Signed Length Quantity]
  \label{def:physics:phyx-mini-0977:phyxminiproblems-problemphyxmini0977-signedlengthquantity}
  \lean{PhyXMiniProblems.ProblemPhyXMini0977.SignedLengthQuantity}
  A signed one-dimensional position or displacement
\end{definition}
\begin{proof}
  This is the declared model definition of Signed Length Quantity.
\end{proof}
\begin{definition}[Mass Quantity]
  \label{def:physics:phyx-mini-0977:phyxminiproblems-problemphyxmini0977-massquantity}
  \lean{PhyXMiniProblems.ProblemPhyXMini0977.MassQuantity}
  A nonnegative, unit-independent mass
\end{definition}
\begin{proof}
  This is the declared model definition of Mass Quantity.
\end{proof}
\begin{definition}[Speed Magnitude Quantity]
  \label{def:physics:phyx-mini-0977:phyxminiproblems-problemphyxmini0977-speedmagnitudequantity}
  \lean{PhyXMiniProblems.ProblemPhyXMini0977.SpeedMagnitudeQuantity}
  \uses{def:physics:phyx-mini-0977:phyxminiproblems-problemphyxmini0977-speeddimension}
  A nonnegative speed magnitude such as the given initial speed
\end{definition}
\begin{proof}
  This is the declared model definition of Speed Magnitude Quantity.
\end{proof}
\begin{definition}[Signed Velocity Quantity]
  \label{def:physics:phyx-mini-0977:phyxminiproblems-problemphyxmini0977-signedvelocityquantity}
  \lean{PhyXMiniProblems.ProblemPhyXMini0977.SignedVelocityQuantity}
  \uses{def:physics:phyx-mini-0977:phyxminiproblems-problemphyxmini0977-speeddimension}
  Signed horizontal velocity, with rightward chosen positive
\end{definition}
\begin{proof}
  This is the declared model definition of Signed Velocity Quantity.
\end{proof}
\begin{definition}[Signed Acceleration Quantity]
  \label{def:physics:phyx-mini-0977:phyxminiproblems-problemphyxmini0977-signedaccelerationquantity}
  \lean{PhyXMiniProblems.ProblemPhyXMini0977.SignedAccelerationQuantity}
  \uses{def:physics:phyx-mini-0977:phyxminiproblems-problemphyxmini0977-accelerationdimension}
  Signed horizontal acceleration, with rightward chosen positive
\end{definition}
\begin{proof}
  This is the declared model definition of Signed Acceleration Quantity.
\end{proof}
\begin{definition}[Magnetic Flux Density Magnitude]
  \label{def:physics:phyx-mini-0977:phyxminiproblems-problemphyxmini0977-magneticfluxdensitymagnitude}
  \lean{PhyXMiniProblems.ProblemPhyXMini0977.MagneticFluxDensityMagnitude}
  \uses{def:physics:phyx-mini-0977:phyxminiproblems-problemphyxmini0977-magneticfluxdensitydimension}
  A nonnegative magnetic-flux-density magnitude
\end{definition}
\begin{proof}
  This is the declared model definition of Magnetic Flux Density Magnitude.
\end{proof}
\begin{definition}[Signed Emf Quantity]
  \label{def:physics:phyx-mini-0977:phyxminiproblems-problemphyxmini0977-signedemfquantity}
  \lean{PhyXMiniProblems.ProblemPhyXMini0977.SignedEmfQuantity}
  \uses{def:physics:phyx-mini-0977:phyxminiproblems-problemphyxmini0977-electromotiveforcedimension}
  Signed motional emf, positive for the counterclockwise loop orientation
\end{definition}
\begin{proof}
  This is the declared model definition of Signed Emf Quantity.
\end{proof}
\begin{definition}[Signed Current Quantity]
  \label{def:physics:phyx-mini-0977:phyxminiproblems-problemphyxmini0977-signedcurrentquantity}
  \lean{PhyXMiniProblems.ProblemPhyXMini0977.SignedCurrentQuantity}
  \uses{def:physics:phyx-mini-0977:phyxminiproblems-problemphyxmini0977-electriccurrentdimension}
  Signed conventional current, positive counterclockwise
\end{definition}
\begin{proof}
  This is the declared model definition of Signed Current Quantity.
\end{proof}
\begin{definition}[Electrical Resistance Quantity]
  \label{def:physics:phyx-mini-0977:phyxminiproblems-problemphyxmini0977-electricalresistancequantity}
  \lean{PhyXMiniProblems.ProblemPhyXMini0977.ElectricalResistanceQuantity}
  \uses{def:physics:phyx-mini-0977:phyxminiproblems-problemphyxmini0977-electricalresistancedimension}
  A nonnegative electrical resistance
\end{definition}
\begin{proof}
  This is the declared model definition of Electrical Resistance Quantity.
\end{proof}
\begin{definition}[Signed Force Quantity]
  \label{def:physics:phyx-mini-0977:phyxminiproblems-problemphyxmini0977-signedforcequantity}
  \lean{PhyXMiniProblems.ProblemPhyXMini0977.SignedForceQuantity}
  \uses{def:physics:phyx-mini-0977:phyxminiproblems-problemphyxmini0977-forcedimension}
  Signed horizontal force, with rightward chosen positive
\end{definition}
\begin{proof}
  This is the declared model definition of Signed Force Quantity.
\end{proof}
\begin{definition}[nonnegative SIReadout]
  \label{def:physics:phyx-mini-0977:phyxminiproblems-problemphyxmini0977-nonnegativesireadout}
  \lean{PhyXMiniProblems.ProblemPhyXMini0977.nonnegativeSIReadout}
  Read a nonnegative dimensionful quantity in coherent SI units
\end{definition}
\begin{proof}
  This is the declared model definition of nonnegative SIReadout.
\end{proof}
\begin{definition}[signed SIReadout]
  \label{def:physics:phyx-mini-0977:phyxminiproblems-problemphyxmini0977-signedsireadout}
  \lean{PhyXMiniProblems.ProblemPhyXMini0977.signedSIReadout}
  Read a signed dimensionful quantity in coherent SI units
\end{definition}
\begin{proof}
  This is the declared model definition of signed SIReadout.
\end{proof}
\begin{definition}[length In Meters]
  \label{def:physics:phyx-mini-0977:phyxminiproblems-problemphyxmini0977-lengthinmeters}
  \lean{PhyXMiniProblems.ProblemPhyXMini0977.lengthInMeters}
  \uses{def:physics:phyx-mini-0977:phyxminiproblems-problemphyxmini0977-lengthquantity, def:physics:phyx-mini-0977:phyxminiproblems-problemphyxmini0977-nonnegativesireadout}
  Read a length in metres
\end{definition}
\begin{proof}
  This is the declared model definition of length In Meters.
\end{proof}
\begin{definition}[signed Length In Meters]
  \label{def:physics:phyx-mini-0977:phyxminiproblems-problemphyxmini0977-signedlengthinmeters}
  \lean{PhyXMiniProblems.ProblemPhyXMini0977.signedLengthInMeters}
  \uses{def:physics:phyx-mini-0977:phyxminiproblems-problemphyxmini0977-signedlengthquantity, def:physics:phyx-mini-0977:phyxminiproblems-problemphyxmini0977-signedsireadout}
  Read a signed position or displacement in metres
\end{definition}
\begin{proof}
  This is the declared model definition of signed Length In Meters.
\end{proof}
\begin{definition}[mass In Kilograms]
  \label{def:physics:phyx-mini-0977:phyxminiproblems-problemphyxmini0977-massinkilograms}
  \lean{PhyXMiniProblems.ProblemPhyXMini0977.massInKilograms}
  \uses{def:physics:phyx-mini-0977:phyxminiproblems-problemphyxmini0977-massquantity, def:physics:phyx-mini-0977:phyxminiproblems-problemphyxmini0977-nonnegativesireadout}
  Read mass in kilograms
\end{definition}
\begin{proof}
  This is the declared model definition of mass In Kilograms.
\end{proof}
\begin{definition}[speed Magnitude In Meters Per Second]
  \label{def:physics:phyx-mini-0977:phyxminiproblems-problemphyxmini0977-speedmagnitudeinmeterspersecond}
  \lean{PhyXMiniProblems.ProblemPhyXMini0977.speedMagnitudeInMetersPerSecond}
  \uses{def:physics:phyx-mini-0977:phyxminiproblems-problemphyxmini0977-speedmagnitudequantity, def:physics:phyx-mini-0977:phyxminiproblems-problemphyxmini0977-nonnegativesireadout}
  Read a speed magnitude in metres per second
\end{definition}
\begin{proof}
  This is the declared model definition of speed Magnitude In Meters Per Second.
\end{proof}
\begin{definition}[signed Velocity In Meters Per Second]
  \label{def:physics:phyx-mini-0977:phyxminiproblems-problemphyxmini0977-signedvelocityinmeterspersecond}
  \lean{PhyXMiniProblems.ProblemPhyXMini0977.signedVelocityInMetersPerSecond}
  \uses{def:physics:phyx-mini-0977:phyxminiproblems-problemphyxmini0977-signedvelocityquantity, def:physics:phyx-mini-0977:phyxminiproblems-problemphyxmini0977-signedsireadout}
  Read signed horizontal velocity in metres per second
\end{definition}
\begin{proof}
  This is the declared model definition of signed Velocity In Meters Per Second.
\end{proof}
\begin{definition}[signed Acceleration In Meters Per Second Squared]
  \label{def:physics:phyx-mini-0977:phyxminiproblems-problemphyxmini0977-signedaccelerationinmeterspersecondsquared}
  \lean{PhyXMiniProblems.ProblemPhyXMini0977.signedAccelerationInMetersPerSecondSquared}
  \uses{def:physics:phyx-mini-0977:phyxminiproblems-problemphyxmini0977-signedaccelerationquantity, def:physics:phyx-mini-0977:phyxminiproblems-problemphyxmini0977-signedsireadout}
  Read signed horizontal acceleration in metres per second squared
\end{definition}
\begin{proof}
  This is the declared model definition of signed Acceleration In Meters Per Second Squared.
\end{proof}
\begin{definition}[magnetic Flux Density In Teslas]
  \label{def:physics:phyx-mini-0977:phyxminiproblems-problemphyxmini0977-magneticfluxdensityinteslas}
  \lean{PhyXMiniProblems.ProblemPhyXMini0977.magneticFluxDensityInTeslas}
  \uses{def:physics:phyx-mini-0977:phyxminiproblems-problemphyxmini0977-magneticfluxdensitymagnitude, def:physics:phyx-mini-0977:phyxminiproblems-problemphyxmini0977-nonnegativesireadout}
  Read magnetic flux density in teslas
\end{definition}
\begin{proof}
  This is the declared model definition of magnetic Flux Density In Teslas.
\end{proof}
\begin{definition}[signed Emf In Volts]
  \label{def:physics:phyx-mini-0977:phyxminiproblems-problemphyxmini0977-signedemfinvolts}
  \lean{PhyXMiniProblems.ProblemPhyXMini0977.signedEmfInVolts}
  \uses{def:physics:phyx-mini-0977:phyxminiproblems-problemphyxmini0977-signedemfquantity, def:physics:phyx-mini-0977:phyxminiproblems-problemphyxmini0977-signedsireadout}
  Read signed emf in volts
\end{definition}
\begin{proof}
  This is the declared model definition of signed Emf In Volts.
\end{proof}
\begin{definition}[signed Current In Amperes]
  \label{def:physics:phyx-mini-0977:phyxminiproblems-problemphyxmini0977-signedcurrentinamperes}
  \lean{PhyXMiniProblems.ProblemPhyXMini0977.signedCurrentInAmperes}
  \uses{def:physics:phyx-mini-0977:phyxminiproblems-problemphyxmini0977-signedcurrentquantity, def:physics:phyx-mini-0977:phyxminiproblems-problemphyxmini0977-signedsireadout}
  Read signed conventional current in amperes
\end{definition}
\begin{proof}
  This is the declared model definition of signed Current In Amperes.
\end{proof}
\begin{definition}[resistance In Ohms]
  \label{def:physics:phyx-mini-0977:phyxminiproblems-problemphyxmini0977-resistanceinohms}
  \lean{PhyXMiniProblems.ProblemPhyXMini0977.resistanceInOhms}
  \uses{def:physics:phyx-mini-0977:phyxminiproblems-problemphyxmini0977-electricalresistancequantity, def:physics:phyx-mini-0977:phyxminiproblems-problemphyxmini0977-nonnegativesireadout}
  Read electrical resistance in ohms
\end{definition}
\begin{proof}
  This is the declared model definition of resistance In Ohms.
\end{proof}
\begin{definition}[signed Force In Newtons]
  \label{def:physics:phyx-mini-0977:phyxminiproblems-problemphyxmini0977-signedforceinnewtons}
  \lean{PhyXMiniProblems.ProblemPhyXMini0977.signedForceInNewtons}
  \uses{def:physics:phyx-mini-0977:phyxminiproblems-problemphyxmini0977-signedforcequantity, def:physics:phyx-mini-0977:phyxminiproblems-problemphyxmini0977-signedsireadout}
  Read signed horizontal force in newtons
\end{definition}
\begin{proof}
  This is the declared model definition of signed Force In Newtons.
\end{proof}
\begin{definition}[Direction Vector]
  \label{def:physics:phyx-mini-0977:phyxminiproblems-problemphyxmini0977-directionvector}
  \lean{PhyXMiniProblems.ProblemPhyXMini0977.DirectionVector}
  A dimensionless direction vector in three-dimensional physical space
\end{definition}
\begin{proof}
  This is the declared model definition of Direction Vector.
\end{proof}
\begin{definition}[into Page Unit Vector]
  \label{def:physics:phyx-mini-0977:phyxminiproblems-problemphyxmini0977-intopageunitvector}
  \lean{PhyXMiniProblems.ProblemPhyXMini0977.intoPageUnitVector}
  \uses{def:physics:phyx-mini-0977:phyxminiproblems-problemphyxmini0977-directionvector}
  Unit direction normal to the diagram and pointing into the page
\end{definition}
\begin{proof}
  This is the declared model definition of into Page Unit Vector.
\end{proof}
\begin{definition}[Diagram Direction]
  \label{def:physics:phyx-mini-0977:phyxminiproblems-problemphyxmini0977-diagramdirection}
  \lean{PhyXMiniProblems.ProblemPhyXMini0977.DiagramDirection}
  Qualitative directions occurring in the raster and force model
\end{definition}
\begin{proof}
  This is the declared model definition of Diagram Direction.
\end{proof}
\begin{definition}[Loop Direction]
  \label{def:physics:phyx-mini-0977:phyxminiproblems-problemphyxmini0977-loopdirection}
  \lean{PhyXMiniProblems.ProblemPhyXMini0977.LoopDirection}
  The two orientations around the rectangular loop
\end{definition}
\begin{proof}
  This is the declared model definition of Loop Direction.
\end{proof}
\begin{definition}[Rail]
  \label{def:physics:phyx-mini-0977:phyxminiproblems-problemphyxmini0977-rail}
  \lean{PhyXMiniProblems.ProblemPhyXMini0977.Rail}
  The two straight horizontal arms of the U-shaped conductor
\end{definition}
\begin{proof}
  This is the declared model definition of Rail.
\end{proof}
\begin{definition}[Figure Object]
  \label{def:physics:phyx-mini-0977:phyxminiproblems-problemphyxmini0977-figureobject}
  \lean{PhyXMiniProblems.ProblemPhyXMini0977.FigureObject}
  Physical and graphical objects visible in the primary raster
\end{definition}
\begin{proof}
  This is the declared model definition of Figure Object.
\end{proof}
\begin{definition}[Figure Label]
  \label{def:physics:phyx-mini-0977:phyxminiproblems-problemphyxmini0977-figurelabel}
  \lean{PhyXMiniProblems.ProblemPhyXMini0977.FigureLabel}
  Symbolic labels printed in 977.png
\end{definition}
\begin{proof}
  This is the declared model definition of Figure Label.
\end{proof}
\begin{definition}[Normal Field Glyph]
  \label{def:physics:phyx-mini-0977:phyxminiproblems-problemphyxmini0977-normalfieldglyph}
  \lean{PhyXMiniProblems.ProblemPhyXMini0977.NormalFieldGlyph}
  Page-normal magnetic-field glyph
\end{definition}
\begin{proof}
  This is the declared model definition of Normal Field Glyph.
\end{proof}
\begin{definition}[Figure Color]
  \label{def:physics:phyx-mini-0977:phyxminiproblems-problemphyxmini0977-figurecolor}
  \lean{PhyXMiniProblems.ProblemPhyXMini0977.FigureColor}
  Colors carrying literal visual information in the supplied image
\end{definition}
\begin{proof}
  This is the declared model definition of Figure Color.
\end{proof}
\begin{definition}[Slidewire Braking Figure]
  \label{def:physics:phyx-mini-0977:phyxminiproblems-problemphyxmini0977-slidewirebrakingfigure}
  \lean{PhyXMiniProblems.ProblemPhyXMini0977.SlidewireBrakingFigure}
  \uses{def:physics:phyx-mini-0977:phyxminiproblems-problemphyxmini0977-diagramdirection, def:physics:phyx-mini-0977:phyxminiproblems-problemphyxmini0977-rail, def:physics:phyx-mini-0977:phyxminiproblems-problemphyxmini0977-figureobject, def:physics:phyx-mini-0977:phyxminiproblems-problemphyxmini0977-figurelabel, def:physics:phyx-mini-0977:phyxminiproblems-problemphyxmini0977-normalfieldglyph, def:physics:phyx-mini-0977:phyxminiproblems-problemphyxmini0977-figurecolor}
  Typed transcription of the non-answer-bearing content of 977.png
\end{definition}
\begin{proof}
  This is the declared model definition of Slidewire Braking Figure.
\end{proof}
\begin{definition}[Slidewire Geometry Model]
  \label{def:physics:phyx-mini-0977:phyxminiproblems-problemphyxmini0977-slidewiregeometrymodel}
  \lean{PhyXMiniProblems.ProblemPhyXMini0977.SlidewireGeometryModel}
  Geometric idealization of the conductor and moving bar
\end{definition}
\begin{proof}
  This is the declared model definition of Slidewire Geometry Model.
\end{proof}
\begin{definition}[Fixed Conductor Model]
  \label{def:physics:phyx-mini-0977:phyxminiproblems-problemphyxmini0977-fixedconductormodel}
  \lean{PhyXMiniProblems.ProblemPhyXMini0977.FixedConductorModel}
  Material idealization used for the U-shaped fixed conductor
\end{definition}
\begin{proof}
  This is the declared model definition of Fixed Conductor Model.
\end{proof}
\begin{definition}[Slidewire Braking Setup]
  \label{def:physics:phyx-mini-0977:phyxminiproblems-problemphyxmini0977-slidewirebrakingsetup}
  \lean{PhyXMiniProblems.ProblemPhyXMini0977.SlidewireBrakingSetup}
  \uses{def:physics:phyx-mini-0977:phyxminiproblems-problemphyxmini0977-lengthquantity, def:physics:phyx-mini-0977:phyxminiproblems-problemphyxmini0977-signedlengthquantity, def:physics:phyx-mini-0977:phyxminiproblems-problemphyxmini0977-massquantity, def:physics:phyx-mini-0977:phyxminiproblems-problemphyxmini0977-speedmagnitudequantity, def:physics:phyx-mini-0977:phyxminiproblems-problemphyxmini0977-signedvelocityquantity, def:physics:phyx-mini-0977:phyxminiproblems-problemphyxmini0977-signedaccelerationquantity, def:physics:phyx-mini-0977:phyxminiproblems-problemphyxmini0977-magneticfluxdensitymagnitude, def:physics:phyx-mini-0977:phyxminiproblems-problemphyxmini0977-signedemfquantity, def:physics:phyx-mini-0977:phyxminiproblems-problemphyxmini0977-signedcurrentquantity, def:physics:phyx-mini-0977:phyxminiproblems-problemphyxmini0977-electricalresistancequantity, def:physics:phyx-mini-0977:phyxminiproblems-problemphyxmini0977-signedforcequantity, def:physics:phyx-mini-0977:phyxminiproblems-problemphyxmini0977-diagramdirection, def:physics:phyx-mini-0977:phyxminiproblems-problemphyxmini0977-loopdirection, def:physics:phyx-mini-0977:phyxminiproblems-problemphyxmini0977-slidewirebrakingfigure, def:physics:phyx-mini-0977:phyxminiproblems-problemphyxmini0977-slidewiregeometrymodel, def:physics:phyx-mini-0977:phyxminiproblems-problemphyxmini0977-fixedconductormodel}
  The apparatus and its independent observables. The real argument of each time-dependent quantity is elapsed time in coherent-SI seconds. Position, velocity, acceleration, emf, current, and force are not defined from one another or from an answer choice
\end{definition}
\begin{proof}
  This is the declared model definition of Slidewire Braking Setup.
\end{proof}
\begin{definition}[Matches Ideal Slidewire Scenario]
  \label{def:physics:phyx-mini-0977:phyxminiproblems-problemphyxmini0977-matchesidealslidewirescenario}
  \lean{PhyXMiniProblems.ProblemPhyXMini0977.MatchesIdealSlidewireScenario}
  \uses{def:physics:phyx-mini-0977:phyxminiproblems-problemphyxmini0977-resistanceinohms, def:physics:phyx-mini-0977:phyxminiproblems-problemphyxmini0977-signedforceinnewtons, def:physics:phyx-mini-0977:phyxminiproblems-problemphyxmini0977-slidewirebrakingsetup}
  The prose-level idealization: the rod is released after receiving its initial speed, mechanical friction vanishes, and the U-shaped fixed conductor's resistance is negligible compared with the rod's resistance (idealized as zero). No stopping-distance formula occurs here
\end{definition}
\begin{proof}
  This is the declared model definition of Matches Ideal Slidewire Scenario.
\end{proof}
\begin{definition}[Matches Primary Slidewire Figure]
  \label{def:physics:phyx-mini-0977:phyxminiproblems-problemphyxmini0977-matchesprimaryslidewirefigure}
  \lean{PhyXMiniProblems.ProblemPhyXMini0977.MatchesPrimarySlidewireFigure}
  \uses{def:physics:phyx-mini-0977:phyxminiproblems-problemphyxmini0977-slidewirebrakingsetup}
  Literal geometry, symbols, directions, and colors read from 977.png
\end{definition}
\begin{proof}
  This is the declared model definition of Matches Primary Slidewire Figure.
\end{proof}
\begin{definition}[Has Physical Slidewire Parameters]
  \label{def:physics:phyx-mini-0977:phyxminiproblems-problemphyxmini0977-hasphysicalslidewireparameters}
  \lean{PhyXMiniProblems.ProblemPhyXMini0977.HasPhysicalSlidewireParameters}
  \uses{def:physics:phyx-mini-0977:phyxminiproblems-problemphyxmini0977-lengthinmeters, def:physics:phyx-mini-0977:phyxminiproblems-problemphyxmini0977-massinkilograms, def:physics:phyx-mini-0977:phyxminiproblems-problemphyxmini0977-speedmagnitudeinmeterspersecond, def:physics:phyx-mini-0977:phyxminiproblems-problemphyxmini0977-magneticfluxdensityinteslas, def:physics:phyx-mini-0977:phyxminiproblems-problemphyxmini0977-resistanceinohms, def:physics:phyx-mini-0977:phyxminiproblems-problemphyxmini0977-slidewirebrakingsetup}
  Positivity and nondegeneracy of all given answer-bearing parameters
\end{definition}
\begin{proof}
  This is the declared model definition of Has Physical Slidewire Parameters.
\end{proof}
\begin{definition}[Models Uniform Into Page Magnetic Field]
  \label{def:physics:phyx-mini-0977:phyxminiproblems-problemphyxmini0977-modelsuniformintopagemagneticfield}
  \lean{PhyXMiniProblems.ProblemPhyXMini0977.ModelsUniformIntoPageMagneticField}
  \uses{def:physics:phyx-mini-0977:phyxminiproblems-problemphyxmini0977-magneticfluxdensityinteslas, def:physics:phyx-mini-0977:phyxminiproblems-problemphyxmini0977-intopageunitvector, def:physics:phyx-mini-0977:phyxminiproblems-problemphyxmini0977-slidewirebrakingsetup}
  The full Physlib magnetic field is uniform in space and time, has magnitude B, and points into the page as indicated by the blue crosses
\end{definition}
\begin{proof}
  This is the declared model definition of Models Uniform Into Page Magnetic Field.
\end{proof}
\begin{definition}[Satisfies Slidewire Kinematics]
  \label{def:physics:phyx-mini-0977:phyxminiproblems-problemphyxmini0977-satisfiesslidewirekinematics}
  \lean{PhyXMiniProblems.ProblemPhyXMini0977.SatisfiesSlidewireKinematics}
  \uses{def:physics:phyx-mini-0977:phyxminiproblems-problemphyxmini0977-signedlengthinmeters, def:physics:phyx-mini-0977:phyxminiproblems-problemphyxmini0977-speedmagnitudeinmeterspersecond, def:physics:phyx-mini-0977:phyxminiproblems-problemphyxmini0977-signedvelocityinmeterspersecond, def:physics:phyx-mini-0977:phyxminiproblems-problemphyxmini0977-signedaccelerationinmeterspersecondsquared, def:physics:phyx-mini-0977:phyxminiproblems-problemphyxmini0977-slidewirebrakingsetup}
  One-dimensional kinematics with the initial conditions from the prose. The position derivative is velocity and the velocity derivative is acceleration
\end{definition}
\begin{proof}
  This is the declared model definition of Satisfies Slidewire Kinematics.
\end{proof}
\begin{definition}[Satisfies Motional Emf Law]
  \label{def:physics:phyx-mini-0977:phyxminiproblems-problemphyxmini0977-satisfiesmotionalemflaw}
  \lean{PhyXMiniProblems.ProblemPhyXMini0977.SatisfiesMotionalEmfLaw}
  \uses{def:physics:phyx-mini-0977:phyxminiproblems-problemphyxmini0977-lengthinmeters, def:physics:phyx-mini-0977:phyxminiproblems-problemphyxmini0977-signedvelocityinmeterspersecond, def:physics:phyx-mini-0977:phyxminiproblems-problemphyxmini0977-magneticfluxdensityinteslas, def:physics:phyx-mini-0977:phyxminiproblems-problemphyxmini0977-signedemfinvolts, def:physics:phyx-mini-0977:phyxminiproblems-problemphyxmini0977-slidewirebrakingsetup}
  Motional induction for a rod of length L moving perpendicular to a uniform field: in the chosen sign convention epsilon = B L v. This is a governing law and does not mention the stopping distance or mass
\end{definition}
\begin{proof}
  This is the declared model definition of Satisfies Motional Emf Law.
\end{proof}
\begin{definition}[Satisfies Closed Loop Ohms Law]
  \label{def:physics:phyx-mini-0977:phyxminiproblems-problemphyxmini0977-satisfiesclosedloopohmslaw}
  \lean{PhyXMiniProblems.ProblemPhyXMini0977.SatisfiesClosedLoopOhmsLaw}
  \uses{def:physics:phyx-mini-0977:phyxminiproblems-problemphyxmini0977-signedemfinvolts, def:physics:phyx-mini-0977:phyxminiproblems-problemphyxmini0977-signedcurrentinamperes, def:physics:phyx-mini-0977:phyxminiproblems-problemphyxmini0977-resistanceinohms, def:physics:phyx-mini-0977:phyxminiproblems-problemphyxmini0977-slidewirebrakingsetup}
  Ohm's law for the complete conducting loop. Keeping the fixed-conductor resistance in this relation makes its later zero-resistance idealization explicit rather than hiding it in a definition of total resistance
\end{definition}
\begin{proof}
  This is the declared model definition of Satisfies Closed Loop Ohms Law.
\end{proof}
\begin{definition}[Satisfies Magnetic Braking Force Law]
  \label{def:physics:phyx-mini-0977:phyxminiproblems-problemphyxmini0977-satisfiesmagneticbrakingforcelaw}
  \lean{PhyXMiniProblems.ProblemPhyXMini0977.SatisfiesMagneticBrakingForceLaw}
  \uses{def:physics:phyx-mini-0977:phyxminiproblems-problemphyxmini0977-lengthinmeters, def:physics:phyx-mini-0977:phyxminiproblems-problemphyxmini0977-signedvelocityinmeterspersecond, def:physics:phyx-mini-0977:phyxminiproblems-problemphyxmini0977-magneticfluxdensityinteslas, def:physics:phyx-mini-0977:phyxminiproblems-problemphyxmini0977-signedcurrentinamperes, def:physics:phyx-mini-0977:phyxminiproblems-problemphyxmini0977-signedforceinnewtons, def:physics:phyx-mini-0977:phyxminiproblems-problemphyxmini0977-slidewirebrakingsetup}
  The magnetic force on the current-carrying slidewire has magnitude I L B and points left, opposing rightward motion. The minus sign refers to the right-positive scalar convention
\end{definition}
\begin{proof}
  This is the declared model definition of Satisfies Magnetic Braking Force Law.
\end{proof}
\begin{definition}[Satisfies Newton Second Law]
  \label{def:physics:phyx-mini-0977:phyxminiproblems-problemphyxmini0977-satisfiesnewtonsecondlaw}
  \lean{PhyXMiniProblems.ProblemPhyXMini0977.SatisfiesNewtonSecondLaw}
  \uses{def:physics:phyx-mini-0977:phyxminiproblems-problemphyxmini0977-massinkilograms, def:physics:phyx-mini-0977:phyxminiproblems-problemphyxmini0977-signedaccelerationinmeterspersecondsquared, def:physics:phyx-mini-0977:phyxminiproblems-problemphyxmini0977-signedforceinnewtons, def:physics:phyx-mini-0977:phyxminiproblems-problemphyxmini0977-slidewirebrakingsetup}
  Newton's second law for the slidewire. Friction and any external drive remain separate force observables; the scenario assumptions set both to zero after release
\end{definition}
\begin{proof}
  This is the declared model definition of Satisfies Newton Second Law.
\end{proof}
\begin{definition}[Has Asymptotic Stopping Distance]
  \label{def:physics:phyx-mini-0977:phyxminiproblems-problemphyxmini0977-hasasymptoticstoppingdistance}
  \lean{PhyXMiniProblems.ProblemPhyXMini0977.HasAsymptoticStoppingDistance}
  \uses{def:physics:phyx-mini-0977:phyxminiproblems-problemphyxmini0977-lengthinmeters, def:physics:phyx-mini-0977:phyxminiproblems-problemphyxmini0977-signedlengthinmeters, def:physics:phyx-mini-0977:phyxminiproblems-problemphyxmini0977-signedvelocityinmeterspersecond, def:physics:phyx-mini-0977:phyxminiproblems-problemphyxmini0977-slidewirebrakingsetup}
  The phrase "distance before coming to rest" is modeled asymptotically: ideal linear magnetic drag makes the speed tend to zero as time tends to infinity. The independent length stoppingDistance is calibrated by the corresponding limit of displacement, with no closed-form value assumed
\end{definition}
\begin{proof}
  This is the declared model definition of Has Asymptotic Stopping Distance.
\end{proof}
\begin{definition}[Answer Choice]
  \label{def:physics:phyx-mini-0977:phyxminiproblems-problemphyxmini0977-answerchoice}
  \lean{PhyXMiniProblems.ProblemPhyXMini0977.AnswerChoice}
  Labels of the four answer choices printed with the problem
\end{definition}
\begin{proof}
  This is the declared model definition of Answer Choice.
\end{proof}
\begin{definition}[displayed Distance In Meters]
  \label{def:physics:phyx-mini-0977:phyxminiproblems-problemphyxmini0977-answerchoice-displayeddistanceinmeters}
  \lean{PhyXMiniProblems.ProblemPhyXMini0977.AnswerChoice.displayedDistanceInMeters}
  \uses{def:physics:phyx-mini-0977:phyxminiproblems-problemphyxmini0977-lengthinmeters, def:physics:phyx-mini-0977:phyxminiproblems-problemphyxmini0977-massinkilograms, def:physics:phyx-mini-0977:phyxminiproblems-problemphyxmini0977-speedmagnitudeinmeterspersecond, def:physics:phyx-mini-0977:phyxminiproblems-problemphyxmini0977-magneticfluxdensityinteslas, def:physics:phyx-mini-0977:phyxminiproblems-problemphyxmini0977-resistanceinohms, def:physics:phyx-mini-0977:phyxminiproblems-problemphyxmini0977-slidewirebrakingsetup, def:physics:phyx-mini-0977:phyxminiproblems-problemphyxmini0977-answerchoice}
  Literal transcription of the four displayed scalar expressions in coherent SI units. This answer metadata does not define the independent stopping- distance observable
\end{definition}
\begin{proof}
  This is the declared model definition of displayed Distance In Meters.
\end{proof}
\begin{definition}[recorded Dataset Answer]
  \label{def:physics:phyx-mini-0977:phyxminiproblems-problemphyxmini0977-recordeddatasetanswer}
  \lean{PhyXMiniProblems.ProblemPhyXMini0977.recordedDatasetAnswer}
  \uses{def:physics:phyx-mini-0977:phyxminiproblems-problemphyxmini0977-answerchoice}
  The answer label recorded in the supplied dataset metadata
\end{definition}
\begin{proof}
  This is the declared model definition of recorded Dataset Answer.
\end{proof}
% --- Archon declaration topology end ---
% --- Archon physics formalization source end ---
